\documentclass{ximera}

\begin{document}
	\author{Stitz-Zeager}
	\xmtitle{Exercises for Exponential Functions}{}

\mfpicnumber{1} \opengraphsfile{ExercisesforExponentialFunctions} % mfpic settings added 


\label{ExercisesforExponentialFunctions}


\begin{question}
In Exercises \ref{graphexpfirsta} - \ref{graphexplasta}, sketch the graph of $g$ by starting with the graph of $f$ and using transformations.  Track at least three points of your choice and the horizontal asymptote through the transformations. State the domain and range of $g$.

\begin{problem}\label{graphexpfirsta}
$f(x) = 2^{x}$, $g(x) = 2^{x} - 1$

\begin{solution}
Use the Desmos graph below with the settings $f(x)=2^x,\,h=0,\,k=-1,\,a=1,\,b=1$

\begin{center} 
\desmos{zr6jvgonqk}{800}{600} 
\end{center}

domain: $(-\infty, \infty)$

range: $(-1, \infty)$

\end{solution}
\end{problem}  

\begin{problem}
$f(x) = \left(\frac{1}{3}\right)^{x}$, $g(x) = \left(\frac{1}{3}\right)^{x-1}$    
\end{problem}

\begin{problem}
$f(x) = 3^{x}$, $g(x) = 3^{-x}+2$

\begin{solution}
Use the Desmos graph below with the settings $f(x)=3^x,\,h=0,\,k=2,\,a=1,\,b=-1$

\begin{center} 
\desmos{jq10j4p3ng}{800}{600} 
\end{center}

domain: $(-\infty, \infty)$

range: $(2, \infty)$

\end{solution}
\end{problem} 

\begin{problem}
$f(x) = 10^{x}$, $g(x) = 10^{\frac{x+1}{2}} - 20$ 
\end{problem}

\begin{problem}
$f(t) = (0.5)^{t}$, $g(t) = 100(0.5)^{0.1t}$

\begin{solution}
Use the Desmos graph below with the settings $f(t)=0.5^t,\,h=0,\,k=0,\,a=100,\,b=0.1$

\begin{center} 
\desmos{sklxycmgqa}{800}{600} 
\end{center}

domain: $(-\infty, \infty)$

range: $(0, \infty)$

\end{solution}
\end{problem}
  
\begin{problem}
$f(t) = (1.25)^{t}$, $g(t) = 1 - (1.25)^{t-2}$
\end{problem}

\begin{problem}
$f(t) = e^{t}$, $g(t) = 8 - e^{-t}$

\begin{solution}
Use the Desmos graph below with the settings $f(t)=e^t,\,h=0,\,k=8,\,a=-1,\,b=-1$

\begin{center} 
\desmos{cawqhkpowm}{800}{600} 
\end{center}

domain: $(-\infty, \infty)$

range: $(-\infty,8)$

\end{solution}
\end{problem}  

\begin{problem}\label{graphexplasta}
$f(t) = e^{t}$, $g(t) = 10e^{-0.1t}$   
\end{problem}  

\end{question}

\begin{question}
In Exercises, \ref{expformfirsta} - \ref{expformlasta}, the graph of an exponential function is given.  Find a formula for the function in the form $F(x) = a \cdot 2^{bx-h}+k$.

\begin{problem}\label{expformfirsta}
Points:  $\left(-2, -\frac{5}{2} \right)$,  $\left(-1, -2 \right)$, $\left(0, -1 \right)$, \\
Asymptote:  $y = -3$. \\

\begin{tikzpicture}[scale=0.6]

% Axes
\draw[->] (-5,0) -- (3,0);
\draw[->] (0,-4) -- (0,6);

% Axis labels
\node at (3,-0.5) {\scriptsize $x$};
\node at (0.5,6) {\scriptsize $y$};

% Tick marks (x)
\foreach \x in {-4,-3,-2,-1,1,2}
  \draw (\x,0.1) -- (\x,-0.1);

% Tick labels (x)
\node at (-4,-0.4) {\scriptsize $-4$};
\node at (-3,-0.4) {\scriptsize $-3$};
\node at (-2,-0.4) {\scriptsize $-2$};
\node at (-1,-0.4) {\scriptsize $-1$};
\node at (1,-0.4) {\scriptsize $1$};
\node at (2,-0.4) {\scriptsize $2$};

% Tick marks (y)
\foreach \y in {-3,-2,-1,1,2,3,4,5}
  \draw (0.1,\y) -- (-0.1,\y);

% Tick labels (y)
\node at (0.4,1) {\scriptsize $1$};
\node at (0.4,2) {\scriptsize $2$};
\node at (0.4,3) {\scriptsize $3$};
\node at (0.4,4) {\scriptsize $4$};
\node at (0.4,5) {\scriptsize $5$};
\node at (0.4,-1) {\scriptsize $-1$};
\node at (0.4,-2) {\scriptsize $-2$};
\node at (0.4,-3) {\scriptsize $-3$};

% Dashed horizontal line y = -3
\draw[dashed] (-5,-3) -- (3,-3);

% Function: y = 2^(x+1) - 3
\draw[domain=-4.5:2.1,smooth,thick,<->] 
  plot (\x,{2^(\x+1) - 3});

% Points
\fill (-2,-2.5) circle (2pt);
\fill (-1,-2) circle (2pt);
\fill (0,-1) circle (2pt);

\end{tikzpicture}

\begin{solution}
$F(x) = 2^{x+1}-3$
\end{solution}

\end{problem} 

\begin{problem}
Points:  $\left(-1, 1 \right)$, $\left(0, 2 \right)$, $\left(1, \frac{5}{2} \right)$, 
Asymptote:  $y = 3$. 

\begin{tikzpicture}[scale=0.6]

% Axes
\draw[->] (-4,0) -- (4,0);
\draw[->] (0,-6) -- (0,4);

% Axis labels
\node at (4,-0.5) {\scriptsize $x$};
\node at (0.5,4) {\scriptsize $y$};

% Tick marks (x)
\foreach \x in {-3,-2,-1,1,2,3}
  \draw (\x,0.1) -- (\x,-0.1);

% Tick labels (x)
\node at (-3,-0.4) {\scriptsize $-3$};
\node at (-2,-0.4) {\scriptsize $-2$};
\node at (-1,-0.4) {\scriptsize $-1$};
\node at (1,-0.4) {\scriptsize $1$};
\node at (2,-0.4) {\scriptsize $2$};
\node at (3,-0.4) {\scriptsize $3$};

% Tick marks (y)
\foreach \y in {-5,-4,-3,-2,-1,1,2,3}
  \draw (0.1,\y) -- (-0.1,\y);

% Tick labels (y)
\node at (0.4,1) {\scriptsize $1$};
\node at (0.4,2) {\scriptsize $2$};
\node at (0.4,3) {\scriptsize $3$};
\node at (0.4,-1) {\scriptsize $-1$};
\node at (0.4,-2) {\scriptsize $-2$};
\node at (0.4,-3) {\scriptsize $-3$};
\node at (0.4,-4) {\scriptsize $-4$};
\node at (0.4,-5) {\scriptsize $-5$};

% Dashed horizontal line y = 3
\draw[dashed] (-4,3) -- (4,3);

% Function: y = 3 - 2^{-x}
\draw[domain=-3:3.1,smooth,thick,<->,samples=100]
  plot (\x,{3 - 2^(-\x)});

% Points
\fill (-1,1) circle (2pt);
\fill (0,2) circle (2pt);
\fill (1,2.5) circle (2pt);

\end{tikzpicture}

\end{problem} 

\begin{problem}
Points:  $\left(\frac{5}{2}, \frac{1}{2} \right)$, $\left(3,1 \right)$, $\left(\frac{7}{2}, 2 \right)$, \\
Asymptote:  $y = 0$. 

\begin{tikzpicture}[scale=0.6]

% Axes
\draw[->] (-1,0) -- (5,0);
\draw[->] (0,-1) -- (0,9);

% Axis labels
\node at (5,-0.5) {\scriptsize $x$};
\node at (0.5,9) {\scriptsize $y$};

% Tick marks (x)
\foreach \x in {1,2,3,4}
  \draw (\x,0.1) -- (\x,-0.1);

% Tick labels (x)
\node at (1,-0.4) {\scriptsize $1$};
\node at (2,-0.4) {\scriptsize $2$};
\node at (3,-0.4) {\scriptsize $3$};
\node at (4,-0.4) {\scriptsize $4$};

% Tick marks (y)
\foreach \y in {1,2,3,4,5,6,7,8}
  \draw (0.1,\y) -- (-0.1,\y);

% Tick labels (y)
\node at (0.4,1) {\scriptsize $1$};
\node at (0.4,2) {\scriptsize $2$};
\node at (0.4,3) {\scriptsize $3$};
\node at (0.4,4) {\scriptsize $4$};
\node at (0.4,5) {\scriptsize $5$};
\node at (0.4,6) {\scriptsize $6$};
\node at (0.4,7) {\scriptsize $7$};
\node at (0.4,8) {\scriptsize $8$};

% Function: y = 2^(2x - 6)
\draw[domain=1.25:4.55,smooth,thick,<->,samples=100]
  plot (\x,{2^(2*\x - 6)});

% Points
\fill (2.5,0.5) circle (2pt);
\fill (3,1) circle (2pt);
\fill (3.5,2) circle (2pt);

\end{tikzpicture}

\begin{solution}
$F(x) = 2^{2x-6}$
\end{solution}

\end{problem}
 


\begin{problem}\label{expformlasta} 
Points:  $\left(-\frac{1}{2}, 6 \right)$, $\left(0,3 \right)$, $\left(\frac{1}{2}, \frac{3}{2} \right)$, \\
Asymptote:  $y = 0$.

\begin{tikzpicture}[scale=0.6]

% Axes
\draw[->] (-4,0) -- (4,0);
\draw[->] (0,-1) -- (0,9);

% Axis labels
\node at (4,-0.5) {\scriptsize $x$};
\node at (0.5,9) {\scriptsize $y$};

% Tick marks (x)
\foreach \x in {-3,-2,-1,1,2,3}
  \draw (\x,0.1) -- (\x,-0.1);

% Tick labels (x)
\node at (-3,-0.4) {\scriptsize $-3$};
\node at (-2,-0.4) {\scriptsize $-2$};
\node at (-1,-0.4) {\scriptsize $-1$};
\node at (1,-0.4) {\scriptsize $1$};
\node at (2,-0.4) {\scriptsize $2$};
\node at (3,-0.4) {\scriptsize $3$};

% Tick marks (y)
\foreach \y in {1,2,3,4,5,6,7,8}
  \draw (0.1,\y) -- (-0.1,\y);

% Tick labels (y) (only selected ones shown, matching mfpic)
\node at (0.4,1) {\scriptsize $1$};
\node at (0.4,2) {\scriptsize $2$};
\node at (0.4,3) {\scriptsize $3$};
\node at (0.4,7) {\scriptsize $7$};
\node at (0.4,8) {\scriptsize $8$};

% Function: y = 3 * 2^(-2x)
\draw[domain=-0.79:2,smooth,thick,<->,samples=100]
  plot (\x,{3*2^(-2*\x)});

% Points
\fill (-0.5,6) circle (2pt);
\fill (0,3) circle (2pt);
\fill (0.5,1.5) circle (2pt);

\end{tikzpicture}
\end{problem}  






\end{question}



\begin{problem}
Find a formula for each graph in Exercises \ref{expformfirsta} - \ref{expformlasta} of the form $G(x) = a \cdot 4^{bx-h} + k$.
  Did you change your solution methodology?    What is the relationship between your answers for $F(x)$ and $G(x)$ for each graph? 

\begin{solution}
Since $2 = 4^{\frac{1}{2}}$, one way to obtain the formulas for $G(x)$ is to use properties of exponents.  For example, $F(x) = 2^{x+1}-3 = \left(4^{\frac{1}{2}}\right)^{x+1} -3 = 4^{\frac{1}{2}(x+1)} - 3 = 4^{\frac{1}{2} x + \frac{1}{2}} - 3$.  In order, the formulas for $G(x)$ are:


\begin{itemize}

\item $G(x) = 4^{\frac{1}{2}x+\frac{1}{2}}-3$

\item  $G(x) = -4^{-\frac{1}{2} x} + 3$

\item $G(x) = 4^{x-3}$

\item  $G(x) =3 \cdot 4^{-x}$

\end{itemize}

\end{solution}
\end{problem}


\begin{problem}\label{morethanoneforexpexercise}
In Example \ref{expfcngraphsex} number \ref{findformulaforexpexample}, we obtained the solution  $F(x) = -2^{x+3} + 4$ as one formula for the given graph by making a simplifying assumption that $a = -1$.  This exercises explores if there are any other solutions for different choices of $a$.

\begin{enumerate}

\item  Show  $G(x) = -4 \cdot 2^{x+1} + 4$ also fits the data for the given graph, and use  properties of exponents to show $G(x) = F(x)$.  (Use the fact that $4 = 2^2$ \ldots)

\item With help from your classmates, find solutions to  Example \ref{expfcngraphsex} number \ref{findformulaforexpexample} using $a = -8$, $a = -16$ and  $a = -\frac{1}{2}$.  Show all your solutions can be rewritten as: $F(x) = -2^{x+3} + 4$.

\item  Using properties of exponents and the fact that the range of $2^{x}$ is $(0, \infty)$, show that any function of the form $f(x) = -a \cdot 2^{bx-h} + k$ for $a> 0$ can be rewritten as $f(x) = - 2^{c} \, 2^{bx-h} + k = -2^{bx-h+c} + k$.  Relabeling, this means every function of the form $f(x) = -a \cdot 2^{bx-h} + k$ with four parameters ($a$, $b$, $h$, and $k$) can be rewritten as $f(x) =  - 2^{bx - H} + k$, a formula with just three parameters:  $b$, $H$, and $k$.  Conclude that  \textit{every} solution to Example \ref{expfcngraphsex} number \ref{findformulaforexpexample} reduces to $F(x) = -2^{x+3} + 4$ .
\end{enumerate}

\end{problem}

\begin{question}
In Exercises \ref{decomposebasicexpfirst} - \ref{decomposebasicexplast}, write the given function as a nontrivial decomposition of functions as directed.

\begin{problem}\label{decomposebasicexpfirst}
For $f(x) = e^{-x} +1 $, find functions $g$ and $h$ so that $f=g+h$.

\begin{solution}
One possible solution is $g(x)  = e^{-x}$ and $h(x) = 1$.
\end{solution}
\end{problem}

\begin{problem}
For $f(x) = e^{2x} - x$, find functions $g$ and $h$ so that $f=g-h$.


\end{problem}  

\begin{problem}
For $f(t) = t^2 e^{-t}$, find functions $g$ and $h$ so that $f=gh$.

\begin{solution}
One possible solution is $g(t) = t^2$ and $h(t) = e^{-t}$.
\end{solution}

\end{problem}

\begin{problem}
For $r(x) = \dfrac{e^{x} - e^{-x}}{e^{x}+e^{-x}}$, find functions $f$ and $g$ so $r = \dfrac{f}{g}$.
\end{problem}

\begin{problem}
For $k(x) = e^{-x^2}$, find functions $f$ and $g$  so that $k = g \circ f$.

\begin{solution}
One possible solution is $f(x) = -x^2$ and $g(x) = e^{x}$.
\end{solution}
\end{problem}

\begin{problem}\label{decomposebasicexplast}
For $s(x) =\sqrt{e^{2x} - 1}$, find functions $f$ and $g$ so $s = g \circ f$.
\end{problem}

\end{question}

\begin{problem}\label{preludetocompoundingexercise}

The amount of money in a savings account, $A(t)$, in dollars,  $t$ years after an initial investment is made is given by: $A(t) = 500(1.05)^{t}$, for $t \geq 0$.

\begin{enumerate}

\item  Find and interpret $A(0)$, $A(1)$, and $A(2)$. 

\begin{solution}
$A(0) = 500$, so the principal is $\$500$.  $A(1) = 525$, so after $1$ year, there is $\$525$ in the savings account.   $A(2) =551.25$, so after $2$ years, there is $\$551.25$ in the savings account. 
\end{solution}

\item  Find and interpret the relative rate of change of $A$ over the intervals $[0,1]$, $[1,2]$, $[0,2]$.

\begin{solution}
The relative rate of change of $A$ over the intervals $[0,1]$ and $[1,2]$ is $0.05$ which means the savings account is growing by $5 \%$ each year for those two years.  Over the interval $[0,2]$, the relative rate of change is $0.1025$ meaning the account has grown by $10.25 \%$ over the course of the first two years.  Note this is greater than the sum of the two rates $5 \% + 5 \% = 10 \%$.  This is due to the `compounding effect' and will be discussed in greater detail in Section \ref{ExpLogApplications}.
\end{solution}

\item  Find, simplify, and interpret the relative rate of change of $A$ over the $[t, t+1]$.  Assume  $t \geq 0$.

\begin{solution}
The relative rate of change of $A$ over the $[t, t+1]$ is $0.05$. This means over the course of one year, the savings account grows by $5 \%$.
\end{solution}

\item Use a graphing utility to estimate how long until the savings account is worth $\$1500$.  Round your answer to the nearest year.

\begin{solution}
Graphing $y= A(t)$ and $y = 1500$, we find they intersect when $t \approx 22.5$ so it takes approximately $22-23$ years for the savings account to grow to $\$1500$ in value.
\end{solution}

\end{enumerate}
  
\end{problem} 

\begin{problem}
Based on census data,\footnote{See \href{http://www.towncharts.com/Ohio/Demographics/Lake-County-OH-Demographics-data.html}{\underline{here}}.} the population of Lake County, Ohio, in 2010 was 230,041 and in 2015, the population was 229,437.  

\begin{enumerate}

\item  Show the percentage change in the population from 2010 to 2015 is approximately $-0.263 \%$.

\item  If this percentage change remains constant, predict the population of Lake County in 2020.

\item  \label{populationfiveyear} Assuming this percentage change per five years remains constant, find an expression for the population $P(t)$ of Lake County where $t$ is the number of five year intervals after 2010.  (So $t = 0$ corresponds to 2010, $t = 1$ corresponds to $2015$, $t = 2$ corresponds to $2020$, etc.)

HINT:  Definitions \ref{expfcnpointbaseform} and \ref{rrc} and ensuing discussion on that page is useful here.

\item  Use your answer to  \ref{populationfiveyear}  to predict the population of Lake County in the year 2017.  

\item  Let $A(t)$  represent the population of Lake County $t$ years after 2010 where the we approximate the percentage change in population per year as $-\frac{0.263 \%}{5} = -0.0526 \%$.    Find a formula for $A(t)$ and compare your predictions with $A(t)$ to those given by $P(t)$.  In particular, what population does each model give for the year 2050?  Discuss any discrepancies with your classmates.

\end{enumerate}
   
\end{problem} 
  
\begin{problem}\label{averageofarc}  
Show that the average rate of change of a function over the interval  $[x, x+2]$ is average of the average rates of change of the function over the intervals $[x,x+1]$ and $[x+1, x+2]$.  Can the same be said for the average rate of change of the function over $[x, x+3]$ and the average of the average rates of change over $[x, x+1]$, $[x+1, x+2]$, and $[x+2, x+3]$?  Generalize.
\end{problem}

\begin{problem}\label{exponentialchangeexercise}  
If $f(x) = b^{x}$ where $b>0$, $b \neq 1$,  show $f(x_{0}+\Delta x) = f(x_{0}) b^{\Delta x}$.
    
\end{problem}

\begin{problem}
Which is larger: $e^{\pi}$ or $\pi^{e}$?  How do you know?  Can you find a proof that doesn't use technology?   
\end{problem}

\begin{problem}\label{derivativeofex}
\begin{enumerate}  \item\label{diffquotex}  Use properties of exponential functions to show that if $f(x) = e^{x}$, then \[ f'(x) = \lim_{h \rightarrow 0} \frac{f(x+h) - f(x)}{h} = \lim_{h \rightarrow 0} e^{x} \, \left( \frac{e^{h} - 1}{h} \right)\]

\item\label{specialexlimit}    Numerically and graphically investigate the limit: $\ds{\lim_{h \rightarrow 0}}$ $\frac{e^{h} - 1}{h}$.

\item  Use parts \ref{diffquotex} and \ref{specialexlimit} to show to your surprise and delight that the derivative of $e^{x}$ is \ldots  $e^{x}$.

\item  Write the equations of the tangent lines to $y = e^{x}$ at the following points:

\begin{enumerate}

\item  $(0,1)$

\item  $(1, e)$

\item $\left(-1, e^{-1} \right)$

\end{enumerate}

\end{enumerate}   
\end{problem}

 


\end{document}
